% Options for packages loaded elsewhere
\PassOptionsToPackage{unicode}{hyperref}
\PassOptionsToPackage{hyphens}{url}
\documentclass[
]{article}
\usepackage{xcolor}
\usepackage[margin=1in]{geometry}
\usepackage{amsmath,amssymb}
\setcounter{secnumdepth}{-\maxdimen} % remove section numbering
\usepackage{iftex}
\ifPDFTeX
  \usepackage[T1]{fontenc}
  \usepackage[utf8]{inputenc}
  \usepackage{textcomp} % provide euro and other symbols
\else % if luatex or xetex
  \usepackage{unicode-math} % this also loads fontspec
  \defaultfontfeatures{Scale=MatchLowercase}
  \defaultfontfeatures[\rmfamily]{Ligatures=TeX,Scale=1}
\fi
\usepackage{lmodern}
\ifPDFTeX\else
  % xetex/luatex font selection
\fi
% Use upquote if available, for straight quotes in verbatim environments
\IfFileExists{upquote.sty}{\usepackage{upquote}}{}
\IfFileExists{microtype.sty}{% use microtype if available
  \usepackage[]{microtype}
  \UseMicrotypeSet[protrusion]{basicmath} % disable protrusion for tt fonts
}{}
\makeatletter
\@ifundefined{KOMAClassName}{% if non-KOMA class
  \IfFileExists{parskip.sty}{%
    \usepackage{parskip}
  }{% else
    \setlength{\parindent}{0pt}
    \setlength{\parskip}{6pt plus 2pt minus 1pt}}
}{% if KOMA class
  \KOMAoptions{parskip=half}}
\makeatother
\usepackage{color}
\usepackage{fancyvrb}
\newcommand{\VerbBar}{|}
\newcommand{\VERB}{\Verb[commandchars=\\\{\}]}
\DefineVerbatimEnvironment{Highlighting}{Verbatim}{commandchars=\\\{\}}
% Add ',fontsize=\small' for more characters per line
\usepackage{framed}
\definecolor{shadecolor}{RGB}{248,248,248}
\newenvironment{Shaded}{\begin{snugshade}}{\end{snugshade}}
\newcommand{\AlertTok}[1]{\textcolor[rgb]{0.94,0.16,0.16}{#1}}
\newcommand{\AnnotationTok}[1]{\textcolor[rgb]{0.56,0.35,0.01}{\textbf{\textit{#1}}}}
\newcommand{\AttributeTok}[1]{\textcolor[rgb]{0.13,0.29,0.53}{#1}}
\newcommand{\BaseNTok}[1]{\textcolor[rgb]{0.00,0.00,0.81}{#1}}
\newcommand{\BuiltInTok}[1]{#1}
\newcommand{\CharTok}[1]{\textcolor[rgb]{0.31,0.60,0.02}{#1}}
\newcommand{\CommentTok}[1]{\textcolor[rgb]{0.56,0.35,0.01}{\textit{#1}}}
\newcommand{\CommentVarTok}[1]{\textcolor[rgb]{0.56,0.35,0.01}{\textbf{\textit{#1}}}}
\newcommand{\ConstantTok}[1]{\textcolor[rgb]{0.56,0.35,0.01}{#1}}
\newcommand{\ControlFlowTok}[1]{\textcolor[rgb]{0.13,0.29,0.53}{\textbf{#1}}}
\newcommand{\DataTypeTok}[1]{\textcolor[rgb]{0.13,0.29,0.53}{#1}}
\newcommand{\DecValTok}[1]{\textcolor[rgb]{0.00,0.00,0.81}{#1}}
\newcommand{\DocumentationTok}[1]{\textcolor[rgb]{0.56,0.35,0.01}{\textbf{\textit{#1}}}}
\newcommand{\ErrorTok}[1]{\textcolor[rgb]{0.64,0.00,0.00}{\textbf{#1}}}
\newcommand{\ExtensionTok}[1]{#1}
\newcommand{\FloatTok}[1]{\textcolor[rgb]{0.00,0.00,0.81}{#1}}
\newcommand{\FunctionTok}[1]{\textcolor[rgb]{0.13,0.29,0.53}{\textbf{#1}}}
\newcommand{\ImportTok}[1]{#1}
\newcommand{\InformationTok}[1]{\textcolor[rgb]{0.56,0.35,0.01}{\textbf{\textit{#1}}}}
\newcommand{\KeywordTok}[1]{\textcolor[rgb]{0.13,0.29,0.53}{\textbf{#1}}}
\newcommand{\NormalTok}[1]{#1}
\newcommand{\OperatorTok}[1]{\textcolor[rgb]{0.81,0.36,0.00}{\textbf{#1}}}
\newcommand{\OtherTok}[1]{\textcolor[rgb]{0.56,0.35,0.01}{#1}}
\newcommand{\PreprocessorTok}[1]{\textcolor[rgb]{0.56,0.35,0.01}{\textit{#1}}}
\newcommand{\RegionMarkerTok}[1]{#1}
\newcommand{\SpecialCharTok}[1]{\textcolor[rgb]{0.81,0.36,0.00}{\textbf{#1}}}
\newcommand{\SpecialStringTok}[1]{\textcolor[rgb]{0.31,0.60,0.02}{#1}}
\newcommand{\StringTok}[1]{\textcolor[rgb]{0.31,0.60,0.02}{#1}}
\newcommand{\VariableTok}[1]{\textcolor[rgb]{0.00,0.00,0.00}{#1}}
\newcommand{\VerbatimStringTok}[1]{\textcolor[rgb]{0.31,0.60,0.02}{#1}}
\newcommand{\WarningTok}[1]{\textcolor[rgb]{0.56,0.35,0.01}{\textbf{\textit{#1}}}}
\usepackage{graphicx}
\makeatletter
\newsavebox\pandoc@box
\newcommand*\pandocbounded[1]{% scales image to fit in text height/width
  \sbox\pandoc@box{#1}%
  \Gscale@div\@tempa{\textheight}{\dimexpr\ht\pandoc@box+\dp\pandoc@box\relax}%
  \Gscale@div\@tempb{\linewidth}{\wd\pandoc@box}%
  \ifdim\@tempb\p@<\@tempa\p@\let\@tempa\@tempb\fi% select the smaller of both
  \ifdim\@tempa\p@<\p@\scalebox{\@tempa}{\usebox\pandoc@box}%
  \else\usebox{\pandoc@box}%
  \fi%
}
% Set default figure placement to htbp
\def\fps@figure{htbp}
\makeatother
\setlength{\emergencystretch}{3em} % prevent overfull lines
\providecommand{\tightlist}{%
  \setlength{\itemsep}{0pt}\setlength{\parskip}{0pt}}
\usepackage{bookmark}
\IfFileExists{xurl.sty}{\usepackage{xurl}}{} % add URL line breaks if available
\urlstyle{same}
\hypersetup{
  hidelinks,
  pdfcreator={LaTeX via pandoc}}

\author{}
\date{\vspace{-2.5em}}

\begin{document}

\section{Mixed Strategies and the Indifference
Principle}\label{sec-mixed-strategies}

\begin{quote}
Why rational players sometimes randomize, how the indifference principle
pins down mixing probabilities, and how to compute and visualize
mixed-strategy Nash equilibria in R.
\end{quote}

\subsection*{Learning objectives}\label{learning-objectives}
\addcontentsline{toc}{subsection}{Learning objectives}

\begin{itemize}
\tightlist
\item
  Explain why pure-strategy equilibria sometimes fail to exist and why
  randomization resolves the problem.
\item
  State and apply the indifference principle: at a mixed Nash
  equilibrium, each player's mixture makes every opponent indifferent
  across the strategies in their support.
\item
  Compute mixed-strategy Nash equilibria for 2x2 games by hand and with
  \texttt{solve\_2x2\_mixed\_nash()}.
\item
  Visualize expected payoffs as a function of mixing probability and
  identify the indifference point graphically.
\end{itemize}

\subsection{Motivation}\label{motivation}

Not every strategic situation has a stable deterministic outcome.
Consider a penalty kick in football. The striker chooses left or right;
the goalkeeper dives left or right. If the goalkeeper could predict the
striker's choice, she would always save the shot. If the striker could
predict the goalkeeper, he would always score. Any predictable pattern
is exploitable, so both players must be unpredictable --- they must
\emph{randomize}.

This is not a peculiarity of sports. Firms randomizing the timing of
sales, tax authorities choosing which returns to audit, and poker
players blending bluffs with genuine bets all face the same logic.
Whenever the interests of the players are sufficiently opposed that no
pure-strategy profile is self-enforcing (see
@ref(sec-nash-equilibrium)), stability requires mixing.

The profound insight of @nash1950 is that this resolution is always
available: every finite game possesses at least one Nash equilibrium
when players are allowed to randomize. This chapter develops the theory
behind that guarantee, introduces the \emph{indifference principle} that
makes mixed equilibria tractable, and implements the computations in R.

\subsection{Theory}\label{theory}

\subsubsection{Mixed strategies defined}\label{mixed-strategies-defined}

In a normal-form game (@ref(sec-normal-form)), a \textbf{pure strategy}
is a definite choice of action. A \textbf{mixed strategy} \(\sigma_i\)
for player \(i\) is a probability distribution over \(i\)'s action set
\(A_i\). We write \(\sigma_i(a_i)\) for the probability that player
\(i\) plays action \(a_i\), with \(\sigma_i(a_i) \geq 0\) for all
\(a_i\) and \(\sum_{a_i \in A_i} \sigma_i(a_i) = 1\).

The \textbf{support} of a mixed strategy is the set of actions played
with positive probability:

\begin{equation}
\text{supp}(\sigma_i) = \{ a_i \in A_i : \sigma_i(a_i) > 0 \}
$$

A pure strategy is a degenerate mixed strategy whose support contains a single action.

### Expected payoffs under mixing

When both players use mixed strategies $\sigma_1$ and $\sigma_2$ in a two-player game, player $i$'s expected payoff is:

$$
U_i(\sigma_1, \sigma_2) = \sum_{a_1 \in A_1} \sum_{a_2 \in A_2} \sigma_1(a_1)\, \sigma_2(a_2)\, u_i(a_1, a_2)
(\#eq:expected-payoff-mixed)
\end{equation}

In a 2x2 game, let \(p\) denote the probability that Player 1 plays
their first action and \(q\) the probability that Player 2 plays their
first action. Then Player 1's expected payoff is:

\begin{equation}
U_1(p, q) = p\,q\,a_{11} + p\,(1-q)\,a_{12} + (1-p)\,q\,a_{21} + (1-p)\,(1-q)\,a_{22}
$$

where $a_{ij}$ are the entries of Player 1's payoff matrix.

### The indifference principle

The key to computing mixed equilibria is a simple but initially counter-intuitive observation.

::: {.rmdnote}
**The Indifference Principle**

At a mixed-strategy Nash equilibrium, each player's mixture must make every other player *indifferent* among all actions in their support. If any action in the support yielded a strictly higher expected payoff, the player would shift all probability to that action, contradicting the assumption of mixing.
:::

This has a striking implication: player $i$'s equilibrium mixing probability is determined not by $i$'s own payoffs, but by the *opponent's* payoffs. Player $i$ mixes in exactly the proportions needed to keep the opponent from exploiting a deviation.

### Computing mixed NE in 2x2 games

Consider a 2x2 game with Player 1's payoff matrix $A$ and Player 2's payoff matrix $B$. Player 2 mixes with probability $q$ on column 1. Player 1 is indifferent between their two rows when:

$$
q \cdot a_{11} + (1-q) \cdot a_{12} = q \cdot a_{21} + (1-q) \cdot a_{22}
$$

Solving for $q$:

$$
q^* = \frac{a_{22} - a_{12}}{a_{11} - a_{21} - a_{12} + a_{22}}
(\#eq:q-star)
\end{equation}

Similarly, Player 1 mixes with probability \(p\) on row 1. Player 2 is
indifferent when:

\begin{equation}
p^* = \frac{b_{22} - b_{12}}{b_{11} - b_{21} - b_{12} + b_{22}}
(\#eq:p-star)
\end{equation}

A mixed NE exists when both \(p^*\) and \(q^*\) fall in \([0, 1]\).

\subsubsection{Nash's existence theorem}\label{nashs-existence-theorem}

\textbf{Theorem: Existence of Nash Equilibrium}

Every finite game --- one with finitely many players and finitely many
actions per player --- has at least one Nash equilibrium in mixed
strategies {[}@nash1950{]}.

The proof relies on Kakutani's fixed-point theorem applied to the
best-response correspondences. The importance of this result cannot be
overstated: it guarantees that for any well-defined finite game, at
least one self-enforcing outcome exists. Without mixed strategies, many
natural games (Matching Pennies, Rock-Paper-Scissors) would have no
equilibrium at all.

\subsection{Implementation in R}\label{implementation-in-r}

\subsubsection{\texorpdfstring{Computing mixed equilibria with
\texttt{solve\_2x2\_mixed\_nash()}}{Computing mixed equilibria with solve\_2x2\_mixed\_nash()}}\label{computing-mixed-equilibria-with-solve_2x2_mixed_nash}

The function \texttt{solve\_2x2\_mixed\_nash()} from
\texttt{R/solvers.R} implements @ref(eq:q-star) and @ref(eq:p-star)
directly.

\begin{Shaded}
\begin{Highlighting}[]
\CommentTok{\# Matching Pennies payoff matrices}
\NormalTok{A\_mp }\OtherTok{\textless{}{-}} \FunctionTok{matrix}\NormalTok{(}\FunctionTok{c}\NormalTok{(}\DecValTok{1}\NormalTok{, }\SpecialCharTok{{-}}\DecValTok{1}\NormalTok{, }\SpecialCharTok{{-}}\DecValTok{1}\NormalTok{, }\DecValTok{1}\NormalTok{), }\AttributeTok{nrow =} \DecValTok{2}\NormalTok{, }\AttributeTok{byrow =} \ConstantTok{TRUE}\NormalTok{,}
               \AttributeTok{dimnames =} \FunctionTok{list}\NormalTok{(}\FunctionTok{c}\NormalTok{(}\StringTok{"Heads"}\NormalTok{, }\StringTok{"Tails"}\NormalTok{), }\FunctionTok{c}\NormalTok{(}\StringTok{"Heads"}\NormalTok{, }\StringTok{"Tails"}\NormalTok{)))}
\NormalTok{B\_mp }\OtherTok{\textless{}{-}} \FunctionTok{matrix}\NormalTok{(}\FunctionTok{c}\NormalTok{(}\SpecialCharTok{{-}}\DecValTok{1}\NormalTok{, }\DecValTok{1}\NormalTok{, }\DecValTok{1}\NormalTok{, }\SpecialCharTok{{-}}\DecValTok{1}\NormalTok{), }\AttributeTok{nrow =} \DecValTok{2}\NormalTok{, }\AttributeTok{byrow =} \ConstantTok{TRUE}\NormalTok{,}
               \AttributeTok{dimnames =} \FunctionTok{list}\NormalTok{(}\FunctionTok{c}\NormalTok{(}\StringTok{"Heads"}\NormalTok{, }\StringTok{"Tails"}\NormalTok{), }\FunctionTok{c}\NormalTok{(}\StringTok{"Heads"}\NormalTok{, }\StringTok{"Tails"}\NormalTok{)))}

\FunctionTok{cat}\NormalTok{(}\StringTok{"Matching Pennies}\SpecialCharTok{\textbackslash{}n}\StringTok{"}\NormalTok{)}
\end{Highlighting}
\end{Shaded}

\begin{verbatim}
#> Matching Pennies
\end{verbatim}

\begin{Shaded}
\begin{Highlighting}[]
\FunctionTok{cat}\NormalTok{(}\StringTok{"Player 1 (Matcher) payoffs:}\SpecialCharTok{\textbackslash{}n}\StringTok{"}\NormalTok{)}
\end{Highlighting}
\end{Shaded}

\begin{verbatim}
#> Player 1 (Matcher) payoffs:
\end{verbatim}

\begin{Shaded}
\begin{Highlighting}[]
\NormalTok{A\_mp}
\end{Highlighting}
\end{Shaded}

\begin{verbatim}
#>       Heads Tails
#> Heads     1    -1
#> Tails    -1     1
\end{verbatim}

\begin{Shaded}
\begin{Highlighting}[]
\FunctionTok{cat}\NormalTok{(}\StringTok{"}\SpecialCharTok{\textbackslash{}n}\StringTok{Player 2 (Mismatcher) payoffs:}\SpecialCharTok{\textbackslash{}n}\StringTok{"}\NormalTok{)}
\end{Highlighting}
\end{Shaded}

\begin{verbatim}
#> 
#> Player 2 (Mismatcher) payoffs:
\end{verbatim}

\begin{Shaded}
\begin{Highlighting}[]
\NormalTok{B\_mp}
\end{Highlighting}
\end{Shaded}

\begin{verbatim}
#>       Heads Tails
#> Heads    -1     1
#> Tails     1    -1
\end{verbatim}

\begin{Shaded}
\begin{Highlighting}[]
\CommentTok{\# Verify no pure{-}strategy NE exists}
\NormalTok{pure\_ne }\OtherTok{\textless{}{-}} \FunctionTok{solve\_2x2\_pure\_nash}\NormalTok{(A\_mp, B\_mp)}
\FunctionTok{cat}\NormalTok{(}\StringTok{"}\SpecialCharTok{\textbackslash{}n}\StringTok{Pure{-}strategy NE:"}\NormalTok{, }\FunctionTok{length}\NormalTok{(pure\_ne), }\StringTok{"}\SpecialCharTok{\textbackslash{}n}\StringTok{"}\NormalTok{)}
\end{Highlighting}
\end{Shaded}

\begin{verbatim}
#> 
#> Pure-strategy NE: 0
\end{verbatim}

\begin{Shaded}
\begin{Highlighting}[]
\CommentTok{\# Compute mixed NE}
\NormalTok{mixed\_ne }\OtherTok{\textless{}{-}} \FunctionTok{solve\_2x2\_mixed\_nash}\NormalTok{(A\_mp, B\_mp)}
\FunctionTok{cat}\NormalTok{(}\FunctionTok{sprintf}\NormalTok{(}\StringTok{"Mixed NE: p = \%.2f (P1 plays Heads), q = \%.2f (P2 plays Heads)}\SpecialCharTok{\textbackslash{}n}\StringTok{"}\NormalTok{,}
\NormalTok{            mixed\_ne}\SpecialCharTok{$}\NormalTok{p, mixed\_ne}\SpecialCharTok{$}\NormalTok{q))}
\end{Highlighting}
\end{Shaded}

\begin{verbatim}
#> Mixed NE: p = 0.50 (P1 plays Heads), q = 0.50 (P2 plays Heads)
\end{verbatim}

\begin{Shaded}
\begin{Highlighting}[]
\CommentTok{\# Expected payoff at mixed NE}
\NormalTok{eu\_mixed }\OtherTok{\textless{}{-}}\NormalTok{ mixed\_ne}\SpecialCharTok{$}\NormalTok{p }\SpecialCharTok{*}\NormalTok{ mixed\_ne}\SpecialCharTok{$}\NormalTok{q }\SpecialCharTok{*}\NormalTok{ A\_mp[}\DecValTok{1}\NormalTok{,}\DecValTok{1}\NormalTok{] }\SpecialCharTok{+}
\NormalTok{            mixed\_ne}\SpecialCharTok{$}\NormalTok{p }\SpecialCharTok{*}\NormalTok{ (}\DecValTok{1} \SpecialCharTok{{-}}\NormalTok{ mixed\_ne}\SpecialCharTok{$}\NormalTok{q) }\SpecialCharTok{*}\NormalTok{ A\_mp[}\DecValTok{1}\NormalTok{,}\DecValTok{2}\NormalTok{] }\SpecialCharTok{+}
\NormalTok{            (}\DecValTok{1} \SpecialCharTok{{-}}\NormalTok{ mixed\_ne}\SpecialCharTok{$}\NormalTok{p) }\SpecialCharTok{*}\NormalTok{ mixed\_ne}\SpecialCharTok{$}\NormalTok{q }\SpecialCharTok{*}\NormalTok{ A\_mp[}\DecValTok{2}\NormalTok{,}\DecValTok{1}\NormalTok{] }\SpecialCharTok{+}
\NormalTok{            (}\DecValTok{1} \SpecialCharTok{{-}}\NormalTok{ mixed\_ne}\SpecialCharTok{$}\NormalTok{p) }\SpecialCharTok{*}\NormalTok{ (}\DecValTok{1} \SpecialCharTok{{-}}\NormalTok{ mixed\_ne}\SpecialCharTok{$}\NormalTok{q) }\SpecialCharTok{*}\NormalTok{ A\_mp[}\DecValTok{2}\NormalTok{,}\DecValTok{2}\NormalTok{]}
\FunctionTok{cat}\NormalTok{(}\FunctionTok{sprintf}\NormalTok{(}\StringTok{"Player 1 expected payoff at mixed NE: \%.2f}\SpecialCharTok{\textbackslash{}n}\StringTok{"}\NormalTok{, eu\_mixed))}
\end{Highlighting}
\end{Shaded}

\begin{verbatim}
#> Player 1 expected payoff at mixed NE: 0.00
\end{verbatim}

\subsubsection{Visualizing the indifference
point}\label{visualizing-the-indifference-point}

The figure below plots Player 1's expected payoff from each pure action
as Player 2 varies \(q\), the probability of playing Heads. At the
indifference point \(q^* = 0.5\), both actions yield the same expected
payoff, so Player 1 is willing to mix.

\begin{Shaded}
\begin{Highlighting}[]
\NormalTok{q\_seq }\OtherTok{\textless{}{-}} \FunctionTok{seq}\NormalTok{(}\DecValTok{0}\NormalTok{, }\DecValTok{1}\NormalTok{, }\AttributeTok{length.out =} \DecValTok{200}\NormalTok{)}

\CommentTok{\# Player 1\textquotesingle{}s expected payoff from Heads: q * 1 + (1{-}q) * ({-}1) = 2q {-} 1}
\CommentTok{\# Player 1\textquotesingle{}s expected payoff from Tails: q * ({-}1) + (1{-}q) * 1 = 1 {-} 2q}
\NormalTok{payoff\_data }\OtherTok{\textless{}{-}} \FunctionTok{tibble}\NormalTok{(}
  \AttributeTok{q =} \FunctionTok{rep}\NormalTok{(q\_seq, }\DecValTok{2}\NormalTok{),}
  \AttributeTok{expected\_payoff =} \FunctionTok{c}\NormalTok{(}\DecValTok{2} \SpecialCharTok{*}\NormalTok{ q\_seq }\SpecialCharTok{{-}} \DecValTok{1}\NormalTok{, }\DecValTok{1} \SpecialCharTok{{-}} \DecValTok{2} \SpecialCharTok{*}\NormalTok{ q\_seq),}
  \AttributeTok{action =} \FunctionTok{rep}\NormalTok{(}\FunctionTok{c}\NormalTok{(}\StringTok{"Heads"}\NormalTok{, }\StringTok{"Tails"}\NormalTok{), }\AttributeTok{each =} \FunctionTok{length}\NormalTok{(q\_seq))}
\NormalTok{)}

\CommentTok{\# Indifference point}
\NormalTok{q\_star }\OtherTok{\textless{}{-}} \FloatTok{0.5}
\NormalTok{eu\_star }\OtherTok{\textless{}{-}} \DecValTok{2} \SpecialCharTok{*}\NormalTok{ q\_star }\SpecialCharTok{{-}} \DecValTok{1}  \CommentTok{\# = 0}

\NormalTok{p\_indiff }\OtherTok{\textless{}{-}} \FunctionTok{ggplot}\NormalTok{(payoff\_data, }\FunctionTok{aes}\NormalTok{(}\AttributeTok{x =}\NormalTok{ q, }\AttributeTok{y =}\NormalTok{ expected\_payoff, }\AttributeTok{colour =}\NormalTok{ action)) }\SpecialCharTok{+}
  \FunctionTok{geom\_line}\NormalTok{(}\AttributeTok{linewidth =} \FloatTok{1.1}\NormalTok{) }\SpecialCharTok{+}
  \FunctionTok{geom\_point}\NormalTok{(}
    \AttributeTok{data =} \FunctionTok{tibble}\NormalTok{(}\AttributeTok{q =}\NormalTok{ q\_star, }\AttributeTok{expected\_payoff =}\NormalTok{ eu\_star),}
    \FunctionTok{aes}\NormalTok{(}\AttributeTok{x =}\NormalTok{ q, }\AttributeTok{y =}\NormalTok{ expected\_payoff),}
    \AttributeTok{colour =}\NormalTok{ okabe\_ito[}\DecValTok{6}\NormalTok{], }\AttributeTok{size =} \DecValTok{4}\NormalTok{, }\AttributeTok{inherit.aes =} \ConstantTok{FALSE}
\NormalTok{  ) }\SpecialCharTok{+}
  \FunctionTok{annotate}\NormalTok{(}\StringTok{"text"}\NormalTok{, }\AttributeTok{x =}\NormalTok{ q\_star }\SpecialCharTok{+} \FloatTok{0.12}\NormalTok{, }\AttributeTok{y =}\NormalTok{ eu\_star }\SpecialCharTok{+} \FloatTok{0.12}\NormalTok{,}
           \AttributeTok{label =} \FunctionTok{paste0}\NormalTok{(}\StringTok{"Indifference point}\SpecialCharTok{\textbackslash{}n}\StringTok{(q* = "}\NormalTok{, q\_star, }\StringTok{")"}\NormalTok{),}
           \AttributeTok{colour =}\NormalTok{ okabe\_ito[}\DecValTok{6}\NormalTok{], }\AttributeTok{size =} \FloatTok{3.5}\NormalTok{, }\AttributeTok{fontface =} \StringTok{"bold"}\NormalTok{) }\SpecialCharTok{+}
  \FunctionTok{geom\_vline}\NormalTok{(}\AttributeTok{xintercept =}\NormalTok{ q\_star, }\AttributeTok{linetype =} \StringTok{"dashed"}\NormalTok{, }\AttributeTok{colour =} \StringTok{"grey50"}\NormalTok{,}
             \AttributeTok{linewidth =} \FloatTok{0.5}\NormalTok{) }\SpecialCharTok{+}
  \FunctionTok{scale\_colour\_manual}\NormalTok{(}
    \AttributeTok{values =} \FunctionTok{c}\NormalTok{(}\StringTok{"Heads"} \OtherTok{=}\NormalTok{ okabe\_ito[}\DecValTok{1}\NormalTok{], }\StringTok{"Tails"} \OtherTok{=}\NormalTok{ okabe\_ito[}\DecValTok{2}\NormalTok{]),}
    \AttributeTok{name =} \StringTok{"Player 1 action"}
\NormalTok{  ) }\SpecialCharTok{+}
  \FunctionTok{scale\_x\_continuous}\NormalTok{(}\AttributeTok{name =} \StringTok{"q  (Player 2\textquotesingle{}s probability of Heads)"}\NormalTok{,}
                     \AttributeTok{breaks =} \FunctionTok{seq}\NormalTok{(}\DecValTok{0}\NormalTok{, }\DecValTok{1}\NormalTok{, }\FloatTok{0.2}\NormalTok{)) }\SpecialCharTok{+}
  \FunctionTok{scale\_y\_continuous}\NormalTok{(}\AttributeTok{name =} \StringTok{"Player 1\textquotesingle{}s expected payoff"}\NormalTok{,}
                     \AttributeTok{breaks =} \FunctionTok{seq}\NormalTok{(}\SpecialCharTok{{-}}\DecValTok{1}\NormalTok{, }\DecValTok{1}\NormalTok{, }\FloatTok{0.5}\NormalTok{)) }\SpecialCharTok{+}
  \FunctionTok{theme\_publication}\NormalTok{() }\SpecialCharTok{+}
  \FunctionTok{labs}\NormalTok{(}\AttributeTok{title =} \StringTok{"Expected Payoffs and the Indifference Principle"}\NormalTok{)}

\NormalTok{p\_indiff}
\end{Highlighting}
\end{Shaded}

\begin{figure}

{\centering \includegraphics[width=0.8\linewidth]{/home/runner/work/strategy-in-r/strategy-in-r/docs/pdf-chapters/04-mixed-strategies_files/figure-latex/indifference-1} 

}

\caption{Player 1's expected payoff from each pure action as a function of Player 2's mixing probability $q$ in Matching Pennies. The lines cross at $q^* = 0.5$, the indifference point where Player 1 is willing to randomize.}\label{fig:indifference}
\end{figure}

\begin{Shaded}
\begin{Highlighting}[]
\FunctionTok{save\_pub\_fig}\NormalTok{(p\_indiff, }\StringTok{"indifference{-}matching{-}pennies"}\NormalTok{)}
\end{Highlighting}
\end{Shaded}

When \(q < 0.5\), Player 2 plays Tails more often, so Player 1 prefers
Tails (the orange line is above). When \(q > 0.5\), Player 1 prefers
Heads (blue line is above). Only at \(q^* = 0.5\) is Player 1 genuinely
indifferent and therefore willing to randomize. The same logic, applied
from Player 2's perspective, pins down \(p^* = 0.5\).

\subsection{Worked example}\label{worked-example}

\subsubsection{Matching Pennies:
step-by-step}\label{matching-pennies-step-by-step}

Two children play Matching Pennies. Each simultaneously shows a coin
Heads or Tails. Player 1 (the Matcher) wins if the coins match; Player 2
(the Mismatcher) wins if they differ. The payoff matrices are as defined
above.

\textbf{Step 1: Check for pure-strategy NE.} For any pure-strategy
profile, one player can profitably deviate. At (Heads, Heads), Player 2
switches to Tails. At (Heads, Tails), Player 1 switches to Tails. No
profile is a mutual best response, so no pure NE exists.

\textbf{Step 2: Apply the indifference principle.} Player 1 is
indifferent between Heads and Tails when:

\[
q \cdot 1 + (1 - q) \cdot (-1) = q \cdot (-1) + (1 - q) \cdot 1
\] \[
2q - 1 = 1 - 2q \implies q^* = \tfrac{1}{2}
\]

Player 2 is indifferent when:

\[
p \cdot (-1) + (1 - p) \cdot 1 = p \cdot 1 + (1 - p) \cdot (-1)
\] \[
1 - 2p = 2p - 1 \implies p^* = \tfrac{1}{2}
\]

\textbf{Step 3: State the equilibrium.} The unique Nash equilibrium is
\((p^*, q^*) = (0.5, 0.5)\). Each player randomizes uniformly. The
expected payoff for both players is 0.

\begin{Shaded}
\begin{Highlighting}[]
\CommentTok{\# Verify with R}
\FunctionTok{cat}\NormalTok{(}\StringTok{"Verification via solve\_2x2\_mixed\_nash():}\SpecialCharTok{\textbackslash{}n}\StringTok{"}\NormalTok{)}
\end{Highlighting}
\end{Shaded}

\begin{verbatim}
#> Verification via solve_2x2_mixed_nash():
\end{verbatim}

\begin{Shaded}
\begin{Highlighting}[]
\NormalTok{result }\OtherTok{\textless{}{-}} \FunctionTok{solve\_2x2\_mixed\_nash}\NormalTok{(A\_mp, B\_mp)}
\FunctionTok{cat}\NormalTok{(}\FunctionTok{sprintf}\NormalTok{(}\StringTok{"  p* = \%.4f, q* = \%.4f}\SpecialCharTok{\textbackslash{}n}\StringTok{"}\NormalTok{, result}\SpecialCharTok{$}\NormalTok{p, result}\SpecialCharTok{$}\NormalTok{q))}
\end{Highlighting}
\end{Shaded}

\begin{verbatim}
#>   p* = 0.5000, q* = 0.5000
\end{verbatim}

\begin{Shaded}
\begin{Highlighting}[]
\CommentTok{\# Full support enumeration}
\NormalTok{all\_eq }\OtherTok{\textless{}{-}} \FunctionTok{support\_enumeration}\NormalTok{(A\_mp, B\_mp)}
\FunctionTok{cat}\NormalTok{(}\FunctionTok{sprintf}\NormalTok{(}\StringTok{"  Pure NE: \%d}\SpecialCharTok{\textbackslash{}n}\StringTok{"}\NormalTok{, }\FunctionTok{length}\NormalTok{(all\_eq}\SpecialCharTok{$}\NormalTok{pure)))}
\end{Highlighting}
\end{Shaded}

\begin{verbatim}
#>   Pure NE: 0
\end{verbatim}

\begin{Shaded}
\begin{Highlighting}[]
\FunctionTok{cat}\NormalTok{(}\FunctionTok{sprintf}\NormalTok{(}\StringTok{"  Mixed NE: p = \%.2f, q = \%.2f}\SpecialCharTok{\textbackslash{}n}\StringTok{"}\NormalTok{, all\_eq}\SpecialCharTok{$}\NormalTok{mixed}\SpecialCharTok{$}\NormalTok{p, all\_eq}\SpecialCharTok{$}\NormalTok{mixed}\SpecialCharTok{$}\NormalTok{q))}
\end{Highlighting}
\end{Shaded}

\begin{verbatim}
#>   Mixed NE: p = 0.50, q = 0.50
\end{verbatim}

\textbf{Step 4: Interpret.} The equilibrium has a natural
interpretation: each player must be maximally unpredictable. Any bias
--- say, playing Heads 60\% of the time --- would be exploited by the
opponent. The mixed NE is the only self-enforcing outcome, consistent
with Nash's existence theorem.

\subsection{Extensions}\label{extensions}

Mixed strategies appear throughout the remainder of this book and in
virtually every branch of applied game theory:

\begin{itemize}
\tightlist
\item
  \textbf{Extensive-form games} (@ref(sec-extensive-form)) use
  behavioral strategies --- independent mixtures at each information set
  --- which are equivalent to mixed strategies under perfect recall.
\item
  In \textbf{Bayesian games} (@ref(sec-bayesian-games)), a player's type
  can be viewed as a private signal that induces mixing from the
  opponent's perspective, blurring the line between mixed and pure
  strategies.
\item
  The \textbf{minimax theorem} of @von-neumann1944 shows that in
  two-player zero-sum games, the mixed NE payoff equals the value of the
  game, and both players can guarantee this value through optimal
  mixing.
\item
  Support enumeration generalizes beyond 2x2 games. For larger games,
  algorithms enumerate all possible support pairs and solve the
  resulting linear systems; see @shoham2009 for computational details.
\end{itemize}

For a thorough treatment of mixed strategies and their properties,
consult @osborne2004 (Chapter 4).

\subsection*{Exercises}\label{exercises}
\addcontentsline{toc}{subsection}{Exercises}

\begin{enumerate}
\def\labelenumi{\arabic{enumi}.}
\item
  \textbf{Asymmetric Matching Pennies.} Consider a variant where Player
  1 receives a payoff of 3 (instead of 1) when both play Heads, while
  all other payoffs remain the same as standard Matching Pennies. Set up
  the payoff matrices, compute the mixed NE by hand using the
  indifference principle, and verify with
  \texttt{solve\_2x2\_mixed\_nash()}. How does the asymmetry affect
  Player 2's mixing probability? Explain the intuition.
\item
  \textbf{Inspection game.} A worker chooses to Shirk or Work; an
  employer chooses to Monitor or Trust. Payoffs: (Work, Trust) yields
  (2, 3), (Work, Monitor) yields (1, 2), (Shirk, Trust) yields (3, 0),
  and (Shirk, Monitor) yields (0, 1). Find the mixed NE. What is the
  equilibrium probability that the employer monitors? Does increasing
  the penalty for caught shirking make shirking less likely in
  equilibrium?
\item
  \textbf{Graphical verification.} Adapt the expected-payoff plot from
  @ref(fig:indifference) to the Battle of the Sexes game (payoffs as in
  @ref(sec-nash-equilibrium)). Plot Player 1's expected payoff from
  Opera and from Football as functions of \(q\). Identify the
  indifference point and verify that it matches the mixed NE from
  \texttt{solve\_2x2\_mixed\_nash()}.
\item
  \textbf{Rock-Paper-Scissors.} Rock-Paper-Scissors is a 3x3 symmetric
  zero-sum game. Argue, using the indifference principle, that the
  unique mixed NE must assign equal probability \(\frac{1}{3}\) to each
  action. (Hint: symmetry means the indifference conditions are
  identical for all pairs of actions.)
\item
  \textbf{Comparative statics.} In a 2x2 game with payoff matrices \(A\)
  and \(B\), suppose we increase \(a_{11}\) (Player 1's payoff when both
  play action 1) while holding all other entries fixed. Using
  @ref(eq:q-star), determine what happens to \(q^*\). Does it increase,
  decrease, or stay the same? Explain why this result is consistent with
  the indifference principle --- and why it surprises many students.
\end{enumerate}

Solutions appear in @ref(sec-solutions).

\end{document}
